% !PRINTED 547
\weberpage{547}

Wenden wir das Theorem \S.\,164, 3. auf nur zwei Factoren an, so ergiebt sich, dass
% [sic] authority prints the Latin u as the exponent; the surrounding formula requires \mu.
\[
(\varepsilon,\alpha)^{\nu}(\varepsilon^{\lambda},\alpha)^{u}
\]
eine Function in $\Omega(\varepsilon)$ ist, wenn $\nu + \lambda\mu \equiv 0 \pmod{m}$. Setzen wir also $\mu = 1$, $\nu = m - \lambda$, so folgt, wenn
\[
\chi_\lambda = b_0^{(\lambda)} + b_1^{(\lambda)}\varepsilon + \cdots + b_{m-1}^{(\lambda)}\varepsilon^{m-1}
\]
eine Grösse in $\Omega(\varepsilon)$ bedeutet,
\[
(\varepsilon,\alpha)^{m-\lambda}(\varepsilon^{\lambda},\alpha) = \chi_\lambda,
\]
also nach (2)
\weq{4}{\sqrt[m]{\psi_\lambda} = \frac{\chi_\lambda}{\left(\sqrt[m]{\psi_1}\right)^{m-\lambda}} = \frac{\chi_\lambda\left(\sqrt[m]{\psi_1}\right)^{\lambda}}{\psi_1},}
und dadurch sind, wenn $\varepsilon$ eine festgehaltene primitive $m^{\mathrm{te}}$ Einheitswurzel bedeutet, die sämmtlichen in (3) vorkommenden Radicale rational durch eines von ihnen, $\sqrt[m]{\psi_1}$, ausgedrückt.

Giebt man diesem einen seine $m$ verschiedenen Werthe, so erhält man aus (3) gerade die $m$ verschiedenen Werthe $\alpha$.

Es ist nur ein Ausnahmefall, in dem dieses Verfahren nicht anwendbar ist, das ist der, wenn $\psi_1 = 0$ ist. Wir können aber durch eine kleine Modification des Verfahrens uns von einem solchen Ausnahmefall frei machen. Dem schicken wir Folgendes voraus.

Es sei $p$ eine in $m$ aufgehende Primzahl und $m = pn$; wie oben sei $\varepsilon$ irgend eine festgehaltene primitive $m^{\mathrm{te}}$ Einheitswurzel. Dann giebt es immer ein durch $p$ nicht theilbares $\lambda$, so dass $(\varepsilon^{\lambda},\alpha)$ von Null verschieden ist. Denn bilden wir nach der Formel \S.\,164, (4) die Differenz $\alpha_n - \alpha_0$, so erhalten wir
\[
m(\alpha_n - \alpha_0) = \sum_{0,\,m-1}^{\lambda}(\varepsilon^{-n\lambda} - 1)(\varepsilon^{\lambda},\alpha).
\]
Nun ist aber $\varepsilon^{-n\lambda} - 1$ immer $= 0$, so oft $\lambda$ durch $p$ theilbar ist, und wenn $(\varepsilon^{\lambda},\alpha)$ in allen anderen Fällen, wo also $\lambda$ nicht durch $p$ theilbar ist, verschwindet, so ist $\alpha_n = \alpha_0$, gegen die Voraussetzung, dass die $\alpha$ alle verschieden sein sollen. Es giebt also wenigstens ein durch $p$ nicht theilbares $\lambda$, so dass $(\varepsilon^{\lambda},\alpha)$ von Null verschieden ist.

Nun zerlegen wir $m$ in seine Primfactoren und setzen
\[
m = p_1 p_2 \dots,
\]
%CONT>>
